\section{Client Heterogeneity}

Clients can vary significantly in their computational resources, networking capabilities and data distributions. These differences affect both what architecture should be searched and how candidate updates should be interpreted.

\subsubsection{Challenge 4: Heterogeneous Deployment Requirements Prevent One Universally Suitable Architecture}

A NAS-in-FL deployment may include devices with different inference-time computation, memory or latency constraints. One architecture sized for the weakest client can underuse stronger hardware, while one sized for the strongest client can be infeasible elsewhere~\cite{fedoras_2022}. The search must therefore produce several target-specific outcomes. Obtaining each through an independent federated search repeats client computation and communication, so the total cost grows with the number of device classes or deployment budgets~\cite{superfednas_2024}.

\noindent\textbf{Solution 4a: Train one supernet for specific deployment constraints.} A weight-sharing supernet can expose subnets with different depths, widths or exits and train them within one federated process. Clients with similar inference constraints can receive one architecture per capability tier~\cite{fedoras_2022,clustered_direct_fednas_2022}, while more finer-grained selection can derive a subnet for each client's hardware limits and validation data~\cite{superfednas_2024,hapfnas_2025,enasfl_2024}. 

This selection determines the deployment outcome, while the capacity matching in Solution 1b determines the workload executed by a client during search. Selection guided by a local predictor can avoid further federated training and ~\cite{superfednas_2024} claims it has reduced training cost by a factor of 11 for 20 targets. Training many subnets together can cause weight-sharing to interfere and leave parameters which are sampled fewer stale~\cite{superfednas_2024}.

\noindent\textbf{Solution 4b: Include resource costs in the architecture search objective.} Supernet-based search strategies can optimize predictive performance subject to measured deployment costs. Latency or multiply-accumulate operation budgets can make outcomes specific to the hardware profile of each client group~\cite{clustered_direct_fednas_2022,load_monitoring_fednas_2025,fgfl_2024,privacy-preserving_nas_across_fl_iot_2021}. The objective can also penalize the number of weights or expected bytes transmitted during subsequent federated training~\cite{multi_objective_evo_fl_2020,fedregnas_2025}. This mechanism changes candidate fitness while retaining the alternatives for search, whereas Solutions 1a and 10a restrict the available modules beforehand. It exposes the trade-off between predictive performance, local inference constraints and recurring model communication before selecting an architecture.

\subsubsection{Challenge 5: Data Heterogeneity Creates Conflicting Architecture Updates}

Architecture updates depend on a given client's distribution. Differences in class balance and feature distributions between clients can cause clients to favor different operations or rank the same candidates differently~\cite{collaborative_nas_for_pfl_2025,superfednas_2024}. Clients that solve related tasks can also require different outputs~\cite{fedpnas_2021}. Clients' architecture updates can also change over time if the client gathers new data~\cite{fedregnas_2025}. 

\noindent\textbf{Solution 5a: Cluster clients by learned architecture preference.} When different data contexts require different architectures, clients can be clustered by the similarity of their locally learned architecture parameters and train one architecture per cluster~\cite{electricity_load_forecasting_fl_darts_2023}. Unlike groups based on latency- or distribution balancing, this mechanism uses intermediate architecture parameters of clients to group clients.

\noindent\textbf{Solution 5b: Optimize the worst-performing subnets across client groups.} A robust supernet search objective can target the candidate with the highest loss on each client partition instead of minimizing only average candidate loss. Sampling the smallest, largest and random subnets provides a good approximation~\cite{superfednas_2024}. By preventing client-subnet combinations with high-loss, client-subnet combinations are protected from being dominated by easier combinations, the objective reduces the effect of interference that depends on the distribution for evaluating shared weights.

\noindent\textbf{Solution 5c: Break the architecture up into shared and task-specific components.} When clients solve related tasks, the search space can separate a shared feature extractor from task-specific output components. Clients can collaborate on the shared structure while retaining locally sampled and fine-tuned components for their prediction task~\cite{fedpnas_2021}.

\noindent\textbf{Cross-cutting solution 6b.} Groups based on resources and represeantive can make the data used to evaluate each candidate approximate the target population~\cite{finch_2024}.

\subsubsection{Challenge 6: Capability-Based Assignment Makes Candidate Evidence Unrepresentative}

Assigning only hardware-feasible subnets allows clients with different capacities to contribute, but it determines which parts of a supernet each client can update. Operations and widths that appear only in larger subnets may receive evidence from a small group of capable clients, while unselected parameters receive no local update~\cite{fedoras_2022,peaches_2024,supernet_trade-off_fednas_2025}. Treating every parameter as though it had been trained uniformly can then distort the shared search state.

Under non-IID data, capability-based assignment also couples architectural capacity with the data of the clients able to execute it. Small candidates may train mainly on low-capability clients and large candidates on well-connected clients~\cite{fedoras_2022}. Candidate quality can then reflect the assignment pattern rather than the architecture alone, even when every selected client returns an update~\cite{hapfnas_2025}. Unequal parameter exposure and data confounding are distinct consequences of the same capability-based assignment mechanism.

\noindent\textbf{Solution 6a: Aggregate parameters according to their actual evaluation time.} The server can aggregate each operation only from clients that updated it and weight those updates by the number of examples that passed through the operation~\cite{fedoras_2022,fgfl_2024,supernet_trade-off_fednas_2025}. The resulting coefficients are parameter-specific and reflect exposure rather than the evaluation unit represented by a complete client update. This prevents untrained parameters from being treated as ordinary client updates and records the evidence available for each part of the supernet.

\noindent\textbf{Solution 6b: Form groups based on resources and represeantive data.} Client groups can be formed under a completion-time or bandwidth constraint and adjusted so that their combined class distribution better represents the target population. The literature implements this by rebalancing bandwidth-based groups according to class distribution or by minimizing each cluster's distance from the population distribution under a completion-time limit~\cite{network_aware_fed_nas_2025,finch_2024}. Applying the mechanism separately to each candidate makes comparisons less dependent on which client links support a particular architecture. Applying it to round-level participant selection can also reduce availability-driven representation bias.

\noindent\textbf{Cross-cutting solution 7d.} Architecture-agnostic distillation transfers predictions or features into candidates that some clients cannot execute, decoupling client knowledge from direct candidate assignment~\cite{hapfnas_2025,collaborative_nas_for_pfl_2025}.
